
We tested whether adding mypy to LLM repair loops improves pass@1 on HumanEval+. It does --- by +8.94 percentage points, consistently across three seeds. We also tested the hypothesized mechanism: that mypy helps specifically for type-error problems. It does not, at least not at k=5 saturation --- both conditions achieve identical 90.9\% repair rates on type-error problems.

The primary contributions of this work are: (1) first controlled evidence that execution+mypy outperforms execution-only repair on HumanEval+ (+8.94pp, three seeds consistent); (2) pre-registered type-specificity refutation (differential=$-0.006$, n=22 type-error problems); (3) benchmark-level mypy signal characterization (70\% on HumanEval+ vs 0\% on MBPP+, all errors in the \texttt{name-defined} category); and (4) Z3 spec generation pipeline with 91.1\% validity rate on arithmetic HumanEval+ problems, establishing feasibility for future formal verification work alongside identification of the current activation-rate (16.1\%) and ceiling-effect limitations.

The mechanism question remains open. General diagnostic context enrichment --- mypy output improving repair quality regardless of error type --- is the most parsimonious interpretation but requires ablation (Condition B-Null) to confirm. If content-specific, the implication is that structured diagnostic text activates latent model knowledge about Python error patterns. If length-driven, the implication is that any additional context improves repair. These two outcomes have substantially different implications for repair loop design.

\textbf{Future directions:} (a) Condition B-Null ablation to distinguish mypy content from prompt length effects (highest priority for mechanism confirmation); (b) per-round type-specificity analysis at k=1..2 before ceiling saturation; (c) multi-model replication (GPT-4o, open-source code models); (d) Z3 CE feedback on harder benchmarks (LiveCodeBench, competition-level problems) with lower base pass@1 and more counterexample headroom; (e) completion of MBPP+ Condition A vs B comparison.

